\documentclass{article}

\usepackage{amsthm, amsmath, amssymb}
\usepackage{tikz-cd}
\usepackage{enumerate}

\newtheorem{theorem}{Theorem}
\newtheorem{lemma}{Lemma}
\newtheorem{proposition}{Proposition}
\newtheorem{definition}{Definition}
\newtheorem{remark}{Remark}

\newcommand{\C}{\mathcal{C}}
\newcommand{\D}{\mathcal{D}}
\newcommand{\K}{\mathcal{K}}
\newcommand{\Set}{\mathbf{Set}}
\newcommand{\Pos}{\mathbf{Pos}}
\newcommand{\id}{\mathrm{id}}
\newcommand{\op}{\mathrm{op}}

\title{Uniqueness of $\K$-ary lax monoidal structures\\ on commutative relative monads}
\author{Andrew Slattery}
\date{June 2026}

\begin{document}

\maketitle

\begin{abstract}
We show that a $\K$-ary lax monoidal structure on a commutative relative monad,
when one exists, is uniquely determined, provided the monad preserves the
cofiltered limits that present infinite products as limits of their finite
subproducts. The proof isolates two independent phenomena: a \emph{reduction}
(the structure at an infinite arity is determined by the finite structures
together with the structure on the constant terminal family), which uses
cofiltered-limit preservation; and the \emph{rigidity of the terminal-family
combinator} (it is forced to be an infinitary meet), which is a purely
equational consequence of associativity and reindexing along surjections, valid
over any cartesian base. No order-enrichment of the codomain is required.
\end{abstract}

\section{Setting}

Let $\C$ and $\D$ be categories with finite products, and let
$J\colon\C\to\D$ be a functor that preserves finite products (a strong monoidal
functor for the cartesian structures). Recall that a \emph{$J$-relative monad}
$(T,\eta,(-)^{*})$ consists of a functor $T\colon\C\to\D$, a unit
$\eta_A\colon JA\to TA$, and a Kleisli extension
$(-)^{*}\colon\D(JA,TB)\to\D(TA,TB)$, subject to the relative monad laws of
Altenkirch--Chapman--Uustalu. We say $T$ is \emph{commutative} when it is
equipped with a commutative tensorial strength; over a base such as $\Set$ this
strength is unique, and it induces the finite lax monoidal structures used
below.

We index products by sets. Fix a class $\K$ of \emph{arities} (sets) such that
$\C$ and $\D$ have $\K$-indexed products, $J$ preserves them, and $\K$ is closed
under finite products, the formation of finite arities, passage to subsets of
its members, and contains the singleton $1$. For an arity $K\in\K$ and a family
$X\colon K\to\C$ we write $X_k$ for its value and $\prod_{k\in K}X_k$ for the
product (formed in $\C$ or in $\D$ as the context demands).

\begin{definition}\label{def:structure}
A \emph{$\K$-ary lax monoidal structure} on the $J$-relative monad $T$ is a
family of morphisms of $\D$, natural in $X\colon K\to\C$,
\[
  \phi_{K,X}\colon \prod_{k\in K} T X_k \longrightarrow T\Bigl(\,\prod_{k\in K} X_k\,\Bigr)
  \qquad (K\in\K),
\]
the inner product formed in $\C$ and the outer in $\D$, satisfying:
\begin{enumerate}[(U)]
  \item[(U)] \emph{Unitality.} $\phi_{1,X}=\id_{TX}$.
  \item[(A)] \emph{Associativity.} For $K\in\K$ and $(J_k)_{k\in K}$ in $\K$
  with $M=\coprod_{k\in K}J_k\in\K$, writing $X\colon M\to\C$ as a doubly-indexed
  family $X_{k,j}$ and $X'_k:=\prod_{j\in J_k}X_{k,j}$:
  \[
  \begin{tikzcd}[column sep=large]
    {\prod_{k}\prod_{j} TX_{k,j}} \arrow[r,"\cong"] \arrow[d,"{\prod_k\phi_{J_k,X_{k,\bullet}}}"'] & {\prod_{m\in M} TX_m} \arrow[dd,"{\phi_{M,X}}"]\\
    {\prod_{k} T X'_k} \arrow[d,"{\phi_{K,X'}}"'] & \\
    {T\bigl(\prod_k X'_k\bigr)} \arrow[r,"\cong"'] & {T\bigl(\prod_{m} X_m\bigr)}
  \end{tikzcd}
  \]
  \item[(R)] \emph{Reindexing along surjections.} For every surjection
  $f\colon K\twoheadrightarrow L$ in $\K$ and every $X\colon L\to\C$, with the
  reindexing maps $f^{*}_{TX}\colon\prod_{l}TX_l\to\prod_{k}TX_{f(k)}$ and
  $f^{*}_{X}\colon\prod_l X_l\to\prod_k X_{f(k)}$ given by precomposition with
  $f$:
  \[
  \begin{tikzcd}[column sep=huge]
    {\prod_{l\in L} TX_l} \arrow[r,"\phi_{L,X}"] \arrow[d,"f^{*}_{TX}"'] & {T(\prod_{l} X_l)} \arrow[d,"T f^{*}_{X}"]\\
    {\prod_{k\in K} TX_{f(k)}} \arrow[r,"\phi_{K,X\circ f}"'] & {T(\prod_{k} X_{f(k)})}
  \end{tikzcd}
  \]
\end{enumerate}
We further require that $\eta$ and the Kleisli extension be monoidal with
respect to $\phi$ (the unit and multiplication axioms). These last two axioms
play no role in what follows and we do not spell them out; the results below use
only \textnormal{(U)}, \textnormal{(A)}, \textnormal{(R)}, naturality, and the
finite uniqueness recorded in Hypothesis~\ref{hyp:fin}.
\end{definition}

\begin{remark}
Reindexing is imposed only along \emph{surjections}. A general reindexing
$f\colon K\to L$ factors as a surjection (a diagonal on products, i.e.
duplication of indices) followed by an injection (a projection, i.e. deletion of
indices); compatibility with the injective part is equivalent to $T$ preserving
the relevant products, which is a strong and often false condition. The
surjective fragment is exactly what is available in the examples of interest and
is exactly what the proof uses.
\end{remark}

We isolate the only input concerning finite arities.

\begin{theorem}[Finite uniqueness; Kock, McDermott--Uustalu]\label{hyp:fin}
For a commutative relative monad over a suitable base \textnormal{(}e.g.
$\C=\Set$, or any base on which tensorial strengths are unique\textnormal{)},
the lax monoidal structure $\phi_{K,X}$ at every \emph{finite} arity $K$ is
uniquely determined.
\end{theorem}

We treat this as known and concentrate on the passage from finite to infinite
arities.

\section{The terminal-family combinator}

Write $M:=T1$, where $1$ is the terminal object of $\C$. For an arity $S\in\K$
let $\mathbf 1\colon S\to\C$ be the constant family at $1$; since
$\prod_{s\in S}1\cong 1$, evaluating the structure at $\mathbf 1$ yields a
morphism of $\D$
\[
  d_S \;:=\; \phi_{S,\mathbf 1}\;\colon\; M^{S}\longrightarrow M .
\]
We call $d_S$ the \emph{combinator} at arity $S$. The whole of this section is a
study of these morphisms; it uses only (U), (A), (R), and is valid over any
cartesian $\D$.

Throughout we reason with generalized elements: a \emph{point of $M$ at stage
$A$} is a morphism $A\to M$ in $\D$, and we compare points by means of the
following relation.

\begin{lemma}\label{lem:semilattice}
Write $m:=d_2\colon M\times M\to M$.
\begin{enumerate}[(i)]
  \item $m$ is commutative, associative, and idempotent.
  \item For each object $A$ of $\D$, the relation on $\D(A,M)$ defined by
  \[
    p\sqsubseteq q \quad:\Longleftrightarrow\quad m\circ\langle p,q\rangle = p
  \]
  is a partial order, natural in $A$, with $m$ computing binary meets:
  $m\circ\langle p,q\rangle=p\wedge q$.
  \item \emph{(Lower bound.)} For every $S\in\K$, every $X$, and each $s\in S$,
  $\;d_S(f)\sqsubseteq f_s\;$ for $f\colon A\to M^{S}$ with components
  $f_s:=\mathrm{pr}_s\circ f$.
  \item \emph{(Constants.)} $d_S\circ\Delta_S=\id_M$, where
  $\Delta_S\colon M\to M^{S}$ is the diagonal; equivalently $d_S$ fixes constant
  families.
  \item \emph{(Interchange / monotonicity.)} $d_2\circ(d_S\times d_S)=d_S\circ
  m^{S}$ up to the symmetry $M^{S}\times M^{S}\cong (M\times M)^{S}$.
  Consequently $d_S$ is monotone: $f\sqsubseteq g$ pointwise implies
  $d_S(f)\sqsubseteq d_S(g)$.
\end{enumerate}
\end{lemma}

\begin{proof}
(i) Commutativity and associativity of $m$ are the instances of (R) along the
transposition $2\to2$ and of (A) along $3=1+1+1$, at the terminal family.
Idempotence is (R) along the fold $\nabla\colon 2\twoheadrightarrow 1$: there
$f^{*}_{TX}\colon M\to M^{2}$ is the diagonal $\Delta$, $f^{*}_X\colon 1\to1$ is
the identity, and $\phi_{1,\mathbf 1}=\id$ by (U), so the square reads
$\id_M = m\circ\Delta$.

(ii) Reflexivity is idempotence; antisymmetry uses commutativity
($p\sqsubseteq q$ and $q\sqsubseteq p$ give $p=m\langle p,q\rangle=m\langle
q,p\rangle=q$); transitivity uses associativity. Naturality in $A$ holds because
$\sqsubseteq$ is defined by an equation between morphisms, which is stable under
precomposition. That $m$ is the meet is precisely the definition of
$\sqsubseteq$ together with $m\langle p,q\rangle\sqsubseteq p,q$, which follows
from (i).

(iii) Apply (A) with outer arity $2$ and blocks $\{s\}$ and $S\setminus s$, so
that $M=\{s\}\sqcup(S\setminus s)=S$. Since $\phi_{\{s\},\mathbf1}=\id$ by (U)
and $\phi_{S\setminus s,\mathbf1}=d_{S\setminus s}$, this gives
$d_S = m\circ(\id\times d_{S\setminus s})$ up to the canonical isomorphism, so
$d_S(f)=m\langle f_s,\,d_{S\setminus s}(f|_{S\setminus s})\rangle\sqsubseteq
f_s$.

(iv) Apply (R) along $\nabla\colon S\twoheadrightarrow 1$ at the terminal family,
exactly as in (i): $f^{*}_{TX}=\Delta_S$, $f^{*}_X=\id$, $\phi_{1,\mathbf1}=\id$,
giving $\id_M=d_S\circ\Delta_S$.

(v) Both sides equal $d_{S\times 2}$ at the terminal family, computed by (A) for
the two factorizations $S\times 2=\coprod_{s\in S}2=\coprod_{i\in 2}S$, identified
by (R) along the symmetry $S\times2\cong 2\times S$. For monotonicity, suppose
$f\sqsubseteq g$ pointwise, i.e. $m\langle f_s,g_s\rangle=f_s$ for all $s$. Then
\[
  m\langle d_S f, d_S g\rangle
  \;\overset{(v)}{=}\;
  d_S\bigl(s\mapsto m\langle f_s,g_s\rangle\bigr)
  \;=\;
  d_S(f),
\]
so $d_S(f)\sqsubseteq d_S(g)$.
\end{proof}

\begin{lemma}[Rigidity of the combinator]\label{lem:combinator}
Let $\phi,\phi'$ be two $\K$-ary lax monoidal structures on $T$ with combinators
$d_S,d'_S$. Then $d_S=d'_S$ for every $S\in\K$. In other words, $d_S$ is the
$S$-ary meet $\bigwedge_{s\in S}$ for the order of
Lemma~\ref{lem:semilattice}, and is in particular unique.
\end{lemma}

\begin{proof}
Both $\phi$ and $\phi'$ induce, by Lemma~\ref{lem:semilattice}(i)--(ii), the
\emph{same} commutative idempotent multiplication $m=d_2=d'_2$ on $M$ (it is the
unique finite structure at arity $2$, by Theorem~\ref{hyp:fin}), hence the same
order $\sqsubseteq$. Fix $S\in\K$ and a point $f\colon A\to M^{S}$. By
Lemma~\ref{lem:semilattice}(iii) for $\phi'$, $d'_S(f)\sqsubseteq f_s$ for all
$s$; equivalently the constant family $\Delta_S\circ d'_S(f)$ satisfies
$\Delta_S\circ d'_S(f)\sqsubseteq f$ pointwise. Applying $d_S$, which is monotone
by Lemma~\ref{lem:semilattice}(v) and fixes constants by (iv),
\[
  d'_S(f)\;=\;d_S\bigl(\Delta_S\circ d'_S(f)\bigr)\;\sqsubseteq\; d_S(f).
\]
By symmetry $d_S(f)\sqsubseteq d'_S(f)$, so antisymmetry gives $d_S(f)=d'_S(f)$.
As $f$ was an arbitrary point (take $A=M^{S}$ and $f=\id$), $d_S=d'_S$.
\end{proof}

\section{From finite to infinite arities}

\begin{theorem}\label{thm:main}
Let $J\colon\C\to\D$ be a finite-product-preserving functor between categories
with finite products, and let $T$ be a commutative $J$-relative monad satisfying
Theorem~\ref{hyp:fin}. Let $K\in\K$ be an arity such that the canonical cone
exhibits
\[
  \prod_{k\in K}X_k \;=\; \varprojlim_{\,J\in\mathcal P_{\mathrm{fin}}(K)}\ \prod_{j\in J}X_j
\]
as a cofiltered limit in $\C$ \emph{(}the transition maps being the
projections\emph{)}, and suppose $T$ preserves this limit. Then any two $\K$-ary
lax monoidal structures on $T$ agree at $K$:
\[
  \phi_{K,X}=\phi'_{K,X}\qquad\text{for every }X\colon K\to\C .
\]
In particular the $\K$-ary lax monoidal structure on $T$, if it exists, is
unique.
\end{theorem}

\begin{proof}
Since $T$ preserves the displayed cofiltered limit, the family
\[
  \bigl\{\,T(\pi_J)\colon T(\textstyle\prod_K X)\to T(\textstyle\prod_J X)\,\bigr\}_{J\in\mathcal P_{\mathrm{fin}}(K)}
\]
is a limit cone in $\D$, hence jointly monic. It therefore suffices to prove
$T(\pi_J)\circ\phi_{K,X}=T(\pi_J)\circ\phi'_{K,X}$ for each finite $J\subseteq
K$.

Fix such a $J$ and let $X^{J}\colon K\to\C$ be the family with $X^{J}_k=X_k$ for
$k\in J$ and $X^{J}_k=1$ otherwise; let $h\colon X\to X^{J}$ be the family
morphism that is the identity on $J$ and the unique map $X_k\to1$ off $J$. Then
$\prod_K X^{J}\cong\prod_J X$ and $\prod_K h=\pi_J$ under this isomorphism, so
naturality of $\phi$ in the family gives
\[
  T(\pi_J)\circ\phi_{K,X}
  \;=\;
  \phi_{K,X^{J}}\circ\prod_{k}T(h_k).
\]
Now apply (A) to the partition $K=J\sqcup(K\setminus J)$, noting
$X^{J}|_{J}=X|_J$ and $X^{J}|_{K\setminus J}=\mathbf 1$:
\[
  \phi_{K,X^{J}}
  \;=\;
  \phi_{2,W}\circ\bigl(\phi_{J,\,X|_J}\times d_{K\setminus J}\bigr),
  \qquad W=\bigl(\textstyle\prod_J X,\;1\bigr),
\]
up to the canonical isomorphisms. In this expression $\phi_{J,X|_J}$ and
$\phi_{2,W}$ are structures at \emph{finite} arities, hence equal for $\phi$ and
$\phi'$ by Theorem~\ref{hyp:fin}; and $d_{K\setminus J}$ is the combinator,
equal for $\phi$ and $\phi'$ by Lemma~\ref{lem:combinator}. Therefore
$T(\pi_J)\circ\phi_{K,X}=T(\pi_J)\circ\phi'_{K,X}$. Joint monicity of the cone
gives $\phi_{K,X}=\phi'_{K,X}$.
\end{proof}

\section{Remarks}

\begin{remark}[What each hypothesis does]
The proof factors into two independent parts, and the hypotheses split
accordingly.
\begin{itemize}
  \item \emph{Commutativity} (Theorem~\ref{hyp:fin}) supplies uniqueness at
  finite arities; it is the only place the monad's strength enters.
  \item \emph{Preservation of the cofiltered limit} $\prod_K=\varprojlim_J\prod_J$
  makes $\{T(\pi_J)\}$ jointly monic and is what reduces the infinite arity $K$
  to its finite truncations together with the single combinator
  $d_{K\setminus J}$ (the proof of Theorem~\ref{thm:main}). It is genuinely used
  here and nowhere else.
  \item \emph{Reindexing along surjections} (R) forces the multiplication
  $d_2$ on $T1$ to be idempotent and the combinators to fix constants
  (Lemma~\ref{lem:semilattice}(i),(iv)); these are what make $d_S$ rigid.
\end{itemize}
\end{remark}

\begin{remark}[No order-enrichment is needed]
The order $\sqsubseteq$ used in Section~2 is the intrinsic semilattice order of
the idempotent commutative multiplication $d_2$ on $T1$; it is automatically a
partial order on each hom-set $\D(A,T1)$, by Lemma~\ref{lem:semilattice}(ii).
At no point is $\D$ assumed to be order-enriched, nor is $T1$ assumed to carry an
order beyond the one its own multiplication defines. The combinator is forced to
be an infinitary meet purely by the equational content of (U), (A), (R). The
codomain $\D$ need only be cartesian.
\end{remark}

\begin{remark}[The role of $\Set\to\Pos$ and the partiality monad]
The motivating instance is $\C=\Set$, $\D=\Pos$, $J$ the discrete-poset
embedding, and $T$ the partiality (lift) monad, for which $T1$ is the
two-element semilattice and $d_2$ is its meet; here $d_S$ is the $S$-ary
conjunction $\forall$, and Theorem~\ref{thm:main} recovers the fact that the
``zip'' on partial elements (defined exactly when all arguments are) is the
unique lax monoidal structure. More generally one may replace the two-element
semilattice by any dominance $\Sigma$, in which case $d_S=\forall_S$ requires
$\forall_{S}$ to preserve $\Sigma$-openness; this is the \emph{compactness}
condition on the arity $S$, and it governs \emph{existence}. Uniqueness, by
contrast, needs no such condition: it is the content of the theorem that the
structure is rigid wherever it exists.
\end{remark}

\begin{remark}[Generality]
The base $\Pos$ plays no part beyond providing the relative structure: the
theorem holds for any finite-product-preserving $J\colon\C\to\D$ between
cartesian categories, with $\C$ carrying the cofiltered limits that present
its $K$-fold products. The unit and multiplication axioms of
Definition~\ref{def:structure} are not used; uniqueness is a property of the
underlying commutative lax monoidal functor, granted finite uniqueness. The
finite-uniqueness input is the natural home for the strength-theoretic results
of Kock and of McDermott--Uustalu.
\end{remark}

\end{document}
